% =====================================================================
% 6 Ensaio 2 — proposta
% Fontes: docs/ensaio2-rota-inversao.md; docs/ars/17 §2 (comentários 1, 2, 6),
% §6 (Hitzig) e §7.6; docs/legislacao/adi-5553-7755-tributacao.md.
% Este capítulo é PLANO, não resultado. O enquadramento é a Decisão D6.
% =====================================================================

\chapter{Ensaio 2: escolha de instrumento}
\label{cap:ensaio2}

\section{Pergunta e ponto de partida}
\label{sec:e2-pergunta}

O Ensaio 2 pergunta se a proibição é o instrumento adequado para a
externalidade da pulverização aérea, em comparação com a taxação do insumo ou
do serviço e com a regulação de desempenho por zonas-tampão. O caso oferece os
três instrumentos em disputa real: a proibição de 2019, a regulação de
desempenho para drones adotada em 2024 e, no plano federal, a desoneração
tributária de agrotóxicos em julgamento no STF (ADI 5553 e ADI 7755), que é a
discussão do instrumento de preço.

A concepção original do ensaio era ancorar a regra de \citet{weitzman1974} na
curva dose-resposta estimada no Ensaio 1. Essa ancoragem não se sustenta, por
quatro razões que o projeto registra.

\begin{alineas}
  \item O $\ATT$ não é o benefício social do controle, e sua derivada não é a
    inclinação do benefício marginal.
  \item A dose do Ensaio 1 mede intensidade de bananicultura e não equivale à
    quantidade que o regulador escolhe, que seria área tratada por via aérea ou
    a exposição resultante.
  \item A proibição é uma cota fixada em zero, isto é, um padrão tecnológico
    \cite{helfand1991}. Ela só coincide com o instrumento de quantidade da
    regra de Weitzman se o ótimo esperado for a solução de canto, e a
    aproximação quadrática local é mais frágil justamente nesse ponto.
  \item O Ensaio 1 não produz, com a precisão obtida, informação utilizável
    sobre a inclinação do benefício marginal.
\end{alineas}

\section{Enquadramento proposto}
\label{sec:e2-enquadramento}

\decisao{6}{enquadrar o Ensaio 2 por suficiência informacional, padrão
tecnológico e objetivo do formulador, em lugar de ancorá-lo no $\ATT$; a
alternativa é adiar o ensaio até haver curvatura identificada.} O ensaio se
organiza em três partes.

\subsection{O que a regra exigiria}

A regra de Weitzman entra como referência para organizar a informação
necessária, e não como fórmula que receba os parâmetros do Ensaio 1. Aplicá-la
pede três objetos definidos sobre a mesma quantidade: a quantidade regulada, a
inclinação do benefício marginal social em relação a ela e a inclinação do
custo marginal agregado de reduzi-la. O ensaio especifica cada objeto para o
caso cearense, indica que dado o mediria e mostra quais deles a evidência
disponível informa. Um exercício de inversão, já implementado, calcula que
combinações de inclinações fariam o ordenamento entre preço e quantidade mudar;
sem calibração, ele não conclui por nenhum instrumento, e esse é o resultado a
reportar.

\subsection{O lado do custo}

A série histórica de custos de produção da banana da Conab traz a operação com
avião como item próprio. Nos sistemas que a utilizam, ela custa de R\$ 280 a
R\$ 520 por hectare e representa cerca de 1,5\% do custo total de produção.
Dos 88 sistemas acompanhados, 4 usam avião e 9 usam trator, e nenhum usa os
dois. Esses dados dão a ordem de
grandeza do custo da proibição. Não dão a curvatura do custo, que depende de
custos marginais dentro do produtor, de custos fixos, de capacidade e escala e
da heterogeneidade entre produtores, nem o custo contrafactual de converter
quem usava avião, porque a escolha de tecnologia é endógena ao sistema de
produção. O ensaio propõe levantar a dispersão do custo de conversão junto ao
setor e à literatura agronômica.

\subsection{O critério de comparação}

A lei declara um objetivo, que é exposição aérea nula. O ensaio propõe comparar
os instrumentos pelo custo e pela confiabilidade com que realizam esse
objetivo sob incerteza, em vez de ordená-los por bem-estar esperado. Sob esse
critério, uma taxa não garante exposição nula quando os custos são incertos, a
zona-tampão depende de fiscalização de desempenho, e a proibição tende a ter custo de
monitoramento menor, por ser observável a presença da aeronave. Essas são
proposições a examinar, e não conclusões. A literatura de
poluição difusa \cite{segerson1988,helfand1995} e a discussão sobre a distância
entre o objetivo do formulador e o do desenho econômico
\cite{hitzig2020normative} fornecem o referencial. A troca de instrumento de
2024, da cota zero para a regulação de desempenho com drones, é o caso a
analisar.

\section{Método e produto esperado}
\label{sec:e2-metodo}

O Ensaio 2 é analítico e documental. Usa o modelo de escolha de instrumento
sob incerteza como estrutura, os dados de custo da Conab como calibração de
ordem de grandeza, os resultados do Ensaio 1 como delimitação do que se sabe
sobre o benefício e os textos legais e as decisões do STF como fonte sobre os
objetivos do formulador. Os fundamentos do acórdão da ADI 6137 ainda serão
lidos na fonte.

O produto esperado é um ensaio que estabeleça (i) que informação a comparação
de instrumentos exige, (ii) que parte dela existe para o caso cearense e (iii)
como os três instrumentos se ordenam sob o critério de custo-efetividade e
confiabilidade, com as hipóteses de cada ordenamento explicitadas. O ensaio não
afirmará qual instrumento maximiza bem-estar esperado, porque a evidência não o
permite.
